% Einheit 2 · Kreisfläche
\setcounter{aufgabe}{10}
\einheitenkopf*{\hypertarget{e2}{}Einheit 2 von 3 · Kreisfläche}
\zweigzeile{Hier lernst du, den Flächeninhalt eines Kreises zu berechnen und aus der Fläche den Radius zu bestimmen · neu oder schon bekannt – je nach Buch · P10 oft · baut auf: Kreisumfang (Einheit 1), Quadrieren und Wurzelziehen, Formeln umstellen}

\begin{aufgabe}{Ich erkenne, welche Formel ich brauche. Kreuze die Formel an, mit der du anfängst. Rechne nicht.}
\begin{geruest}
\gz{a}{gegeben: Radius, gesucht: Flächeninhalt}{\kreuz{$u = \pi \cdot d$}\quad\kreuz{$A = \pi \cdot r^2$}}
\gz{b}{gegeben: Durchmesser, gesucht: Umfang}{\kreuz{$u = \pi \cdot d$}\quad\kreuz{$A = \pi \cdot r^2$}}
\gz{c}{gegeben: Durchmesser, gesucht: Flächeninhalt}{\kreuz{$u = \pi \cdot d$}\quad\kreuz{$A = \pi \cdot r^2$}}
\gz{d}{gegeben: Umfang, gesucht: Radius}{\kreuz{$u = \pi \cdot d$}\quad\kreuz{$A = \pi \cdot r^2$}}
\gz{e}{gegeben: Flächeninhalt, gesucht: Radius}{\kreuz{$u = \pi \cdot d$}\quad\kreuz{$A = \pi \cdot r^2$}}
\end{geruest}
\end{aufgabe}

\begin{aufgabe}{Ich kann den Flächeninhalt eines Kreises berechnen. Rechne mit der $\pi$-Taste und runde auf eine Stelle nach dem Komma.}
\beispiel{r = 6\,\text{cm} \quad\rightarrow\quad A &= \pi \cdot 6^2\,\text{cm}^2 \approx 113{,}1\,\text{cm}^2}
Berechne den Flächeninhalt: $A = \pi \cdot r^2$.
\begin{teilezwei}
\tz $r = 1$\,cm \quad \feld[cm²]{A} & \tz $r = 7$\,cm \quad \feld[cm²]{A} \\
\tz $r = 9$\,m \quad \feld[m²]{A} & \tz $r = 12$\,cm \quad \feld[cm²]{A} \\
\end{teilezwei}
Ist der Durchmesser gegeben, halbierst du ihn zuerst. Berechne den Flächeninhalt.
\begin{teilezwei}
\tz $d = 16$\,cm \quad \feld[cm²]{A} & \tz $r = 2{,}5$\,cm \quad \feld[cm²]{A} \\
\end{teilezwei}
\end{aufgabe}

\begin{aufgabe}{Ich kann den Flächeninhalt von Halbkreis und Viertelkreis berechnen. Runde auf eine Stelle nach dem Komma.}
\begin{kreis}[1.3]\mittelpunkt{M}\sektor{0}{180}{}\radius{0}{r}\end{kreis}\hspace{1.5cm}\begin{kreis}[1.3]\mittelpunkt{M}\sektor{0}{90}{}\radius{0}{r}\end{kreis}

\begin{teile}
\teil Halbkreis mit $r = 8$\,cm \quad \feld[cm²]{A}
\teil Viertelkreis mit $r = 11$\,cm \quad \feld[cm²]{A}
\teil Halbkreis mit dem Durchmesser $d = 9$\,m \quad \feld[m²]{A}
\end{teile}
Schreibe einen Term oder benenne die Figur.
\begin{teile}
\teil Schreibe einen Term für den Flächeninhalt eines Viertelkreises mit dem Radius $r$.
\teil Welcher Teil eines Kreises hat den Flächeninhalt $\frac34 \cdot \pi \cdot r^2$?
\end{teile}
\end{aufgabe}

\begin{aufgabe}{Ich kann aus einer Größe eines Kreises die anderen berechnen. Ergänze die Tabelle. Runde auf eine Stelle nach dem Komma und rechne mit den ungerundeten Werten weiter.}
\sachtabelle{lcccc}{ & Radius $r$ & Durchmesser $d$ & Umfang $u$ & Flächeninhalt $A$}{a) & 1,5\,m & \leerzelle & \leerzelle & \leerzelle \\ b) & \leerzelle & 30\,cm & \leerzelle & \leerzelle \\ c) & \leerzelle & \leerzelle & 40\,cm & \leerzelle}
\end{aufgabe}

\begin{aufgabe}{Ich kann aus dem Flächeninhalt den Radius berechnen. Runde auf eine Stelle nach dem Komma.}
\beispiel{A = 50\,\text{cm}^2 \quad\rightarrow\quad r &= \sqrt{50 : \pi}\,\text{cm} \approx 4{,}0\,\text{cm}}
Berechne den Radius.
\begin{teilezwei}
\tz $A = 30$\,cm² \quad \feld[cm]{r} & \tz $A = 200$\,m² \quad \feld[m]{r} \\
\end{teilezwei}
Berechne den Durchmesser.
\begin{teilezwei}
\tz $A = 1000$\,cm² \quad \feld[cm]{d} & \\
\end{teilezwei}
\begin{teile}
\teil Ein rundes Trampolin hat einen Durchmesser von 3,66\,m. Berechne den Flächeninhalt der Sprungfläche und runde auf zwei Stellen nach dem Komma. (P10 2016)
\end{teile}
\end{aufgabe}

\begin{aufgabe}{Ich finde den Fehler bei der Kreisfläche. Mia berechnet den Flächeninhalt eines Kreises mit dem Radius $r = 17$\,cm so:}
\rechnung{A &= \pi \cdot r^2 \\ A &= \pi \cdot 2 \cdot 17\,\text{cm}^2 \\ A &\approx 106{,}8\,\text{cm}^2}
\begin{teile}
\teil Finde den Fehler, benenne ihn und korrigiere die Rechnung.
\teil Berechne den Flächeninhalt eines Kreises mit dem Radius $r = 0{,}8$\,m.
\end{teile}
\end{aufgabe}

\begin{aufgabe}{Ich kann unterscheiden, ob der Umfang oder der Flächeninhalt gesucht ist. Entscheide, was gesucht ist, und berechne. Runde auf eine Stelle nach dem Komma.}
\begin{teile}
\teil Ein runder Spiegel hat den Radius 25\,cm. Wie viel Glas braucht man für den Spiegel?
\teil Um den Spiegel wird eine Holzleiste geklebt. Wie lang muss die Leiste sein?
\teil Ein runder Teich hat einen Durchmesser von 4,2\,m. Wie lang wird ein Zaun, der direkt am Rand entlangläuft?
\teil Der Boden des Teichs wird mit Folie ausgelegt. Wie viel Folie braucht man mindestens?
\end{teile}
\end{aufgabe}

\begin{aufgabe}{Ich kann begründen, warum die Kreisfläche $\pi \cdot r^2$ ist. Ein Kreis mit dem Radius $r$ wird in 12 gleiche Stücke geschnitten. Die Stücke werden abwechselnd nebeneinandergelegt.}
\begin{minipage}[c]{4.5cm}\tortenkreis\end{minipage}\hspace{1cm}\begin{minipage}[c]{8cm}\tortenrechteck\end{minipage}

\begin{teile}
\teil Wie lang ist die untere Seite der neuen Figur ungefähr? Begründe mit dem Umfang des Kreises.
\teil Wie hoch ist die neue Figur ungefähr?
\teil Begründe damit, dass der Kreis den Flächeninhalt $A = \pi \cdot r^2$ hat.
\teil Tim sagt: „Wenn ich den Radius verdopple, verdoppelt sich auch die Fläche.“ Stimmt das? Begründe.
\end{teile}
\end{aufgabe}

\begin{aufgabe}{Ich kann mit Kreisflächen Preise vergleichen. Eine Pizzeria verkauft eine kleine Pizza mit 26\,cm Durchmesser für 8,50\,€ und eine große Pizza mit 36\,cm Durchmesser für 14,00\,€.}
\begin{teile}
\teil Berechne den Flächeninhalt beider Pizzen.
\teil Berechne für beide Pizzen, was 100\,cm² Pizza kosten. Runde auf Cent.
\teil Welche Pizza ist im Verhältnis zu ihrer Größe günstiger? Begründe mit deinen Ergebnissen.
\end{teile}
\end{aufgabe}
